%% Beginning of file 'sample7.tex'
%%
%% Version 7. Created January 2025.  
%%
%% AASTeX v7 calls the following external packages:
%% times, hyperref, ifthen, hyphens, longtable, xcolor, 
%% bookmarks, array, rotating, ulem, and lineno 
%%
%% RevTeX is no longer used in AASTeX v7.
%%
\documentclass[linenumbers,twocolumn]{aastex7}

\usepackage{savesym}
\usepackage{graphicx}
\usepackage{amsmath}
\usepackage{bm}
\savesymbol{tablenum}
\usepackage{siunitx}
\restoresymbol{SIX}{tablenum}
\usepackage{xcolor}
\usepackage{listings}
\usepackage{multirow}
\usepackage{tikz}

\lstset{%
basicstyle=\ttfamily,
breaklines = false,
}

\sisetup{range-phrase=\,--\,} 
\sisetup{range-units=single} 
\DeclareSIUnit\erg{erg}
\DeclareSIUnit\pixel{pixel}
%%
%% This initial command takes arguments that can be used to easily modify 
%% the output of the compiled manuscript. Any combination of arguments can be 
%% invoked like this:
%%
%% \documentclass[argument1,argument2,argument3,...]{aastex7}
%%
%% Six of the arguments are typestting options. They are:
%%
%%  twocolumn   : two text columns, 10 point font, single spaced article.
%%                This is the most compact and represent the final published
%%                derived PDF copy of the accepted manuscript from the publisher
%%  default     : one text column, 10 point font, single spaced (default).
%%  manuscript  : one text column, 12 point font, double spaced article.
%%  preprint    : one text column, 12 point font, single spaced article.  
%%  preprint2   : two text columns, 12 point font, single spaced article.
%%  modern      : a stylish, single text column, 12 point font, article with
%% 		  wider left and right margins. This uses the Daniel
%% 		  Foreman-Mackey and David Hogg design.
%%
%% Note that you can submit to the AAS Journals in any of these 6 styles.
%%
%% There are other optional arguments one can invoke to allow other stylistic
%% actions. The available options are:
%%
%%   astrosymb    : Loads Astrosymb font and define \astrocommands. 
%%   tighten      : Makes baselineskip slightly smaller, only works with 
%%                  the twocolumn substyle.
%%   times        : uses times font instead of the default.
%%   linenumbers  : turn on linenumbering. Note this is mandatory for AAS
%%                  Journal submissions and revisions.
%%   trackchanges : Shows added text in bold.
%%   longauthor   : Do not use the more compressed footnote style (default) for 
%%                  the author/collaboration/affiliations. Instead print all
%%                  affiliation information after each name. Creates a much 
%%                  longer author list but may be desirable for short 
%%                  author papers.
%% twocolappendix : make 2 column appendix.
%%   anonymous    : Do not show the authors, affiliations, acknowledgments,
%%                  and author contributions for dual anonymous review.
%%  resetfootnote : Reset footnotes to 1 in the body of the manuscript.
%%                  Useful when there are a lot of authors and affiliations
%%		    in the front matter.
%%   longbib      : Print article titles in the references. This option
%% 		    is mandatory for PSJ manuscripts.
%%
%% Since v6, AASTeX has included \hyperref support. While we have built in 
%% specific %% defaults into the classfile you can manually override them 
%% with the \hypersetup command. For example,
%%
%% \hypersetup{linkcolor=red,citecolor=green,filecolor=cyan,urlcolor=magenta}
%%
%% will change the color of the internal links to red, the links to the
%% bibliography to green, the file links to cyan, and the external links to
%% magenta. Additional information on \hyperref options can be found here:
%% https://www.tug.org/applications/hyperref/manual.html#x1-40003
%%
%% The "bookmarks" has been changed to "true" in hyperref
%% to improve the accessibility of the compiled pdf file.
%%
%% If you want to create your own macros, you can do so
%% using \newcommand. Your macros should appear before
%% the \begin{document} command.
%%
\newcommand{\vdag}{(v)^\dagger}
\newcommand\aastex{AAS\TeX}
\newcommand\latex{La\TeX}
\DeclareMathOperator\erf{erf}
\DeclareMathOperator*{\argmax}{arg\,max}
\DeclareMathOperator*{\argmin}{arg\,min}
%%%%%%%%%%%%%%%%%%%%%%%%%%%%%%%%%%%%%%%%%%%%%%%%%%%%%%%%%%%%%%%%%%%%%%%%%%%%%%%%
%%
%% The following section outlines numerous optional output that
%% can be displayed in the front matter or as running meta-data.
%%
%% Running header information. A short title on odd pages and 
%% short author list on even pages. Note that this
%% information may be modified in production.
%%\shorttitle{AASTeX v7 Sample article}
%%\shortauthors{The Terra Mater collaboration}
%%
%% Include dates for submitted, revised, and accepted.
%%\received{February 1, 2025}
%%\revised{March 1, 2025}
%%\accepted{\today}
%%
%% Indicate AAS Journal the manuscript was submitted to.
% \submitjournal{APJ}
%% Note that this command adds "Submitted to " the argument.
%%
%% You can add a light gray and diagonal water-mark to the first page 
%% with this command:
% \watermark{text}
%% where "text", e.g. DRAFT, is the text to appear.  If the text is 
%% long you can control the water-mark size with:
%% \setwatermarkfontsize{dimension}
%% where dimension is any recognized LaTeX dimension, e.g. pt, in, etc.
%%%%%%%%%%%%%%%%%%%%%%%%%%%%%%%%%%%%%%%%%%%%%%%%%%%%%%%%%%%%%%%%%%%%%%%%%%%%%%%%
%%
%% Use this command to indicate a subdirectory where figures are located.
%%\graphicspath{{./}{figures/}}
%% This is the end of the preamble.  Indicate the beginning of the
%% manuscript itself with \begin{document}.

% \graphicspath{{./}{figures/}}

\begin{document}

\title{Solar Slitless Spectroscopy with Learned Reconstruction}

%% A significant change from AASTeX v6+ is in the author blocks. Now an email
%% address is required for each author. This means that each author requires
%% at least one of the following:
%%
%% \author
%% \affiliation
%% \email
%%
%% If these three commands are not available for each author, the latex
%% compiler will issue an error and if you force the latex compiler to continue,
%% it will generate an incomplete pdf.
%%
%% Multiple \affiliation commands are allowed and authors can also include
%% an optional \altaffiliation to indicate a status, i.e. Hubble Fellow. 
%% while affiliations are indexed as footnotes, altaffiliations are noted with
%% with a non-numeric footnote that is set away from the numeric \affiliation 
%% footnotes. NOTE that if an \altaffiliation command is used it must 
%% come BEFORE the \affiliation call, right after the \author command, in 
%% order to place the footnotes in the proper location. Because non-numeric
%% symbols are used, \altaffiliation should be used sparingly.
%%
%% In v7 the \author command takes an optional argument which provides 
%% additional metadata about the author. Authors can provide the 16 digit 
%% ORCID, the surname (family or last) name, the given (first or fore-) name, 
%% and a name suffix, e.g. "Jr.". The syntax is:
%%
%% \author[orcid=0000-0002-9072-1121,gname=Gregory,sname=Schwarz]{Greg Schwarz}
%%
%% This name metadata in not shown, it is only for parsing by the peer review
%% system so authors can be more easily identified. This name information will
%% also be sent to the publisher so they can include it in the CROSSREF 
%% metadata. Including an orcid will hyperlink the author name to the 
%% author's ORCID page. Note that  during compilation, LaTeX will do some 
%% limited checking of the format of the ID to make sure it is valid. If 
%% the "orcid-ID.png" image file is  present or in the LaTeX pathway, the 
%% ORCID icon will appear next to the authors name.
%%
%% Even though emails are now required for each author, the \email does not
%% produce output in the compiled manuscript unless the optional "show" command
%% is used. For example,
%%
%% \email[show]{greg.schwarz@aas.org}
%%
%% All "shown" emails are show in the bottom left of the first page. Due to
%% space constraints, only a few emails should be shown. 
%%
%% To identify a corresponding author, use the \correspondingauthor command.
%% The command appends "Corresponding Author: " to the argument it appears at
%% the bottom left of the first page like the output from \email. 


\author[orcid=0009-0009-4606-6842, gname='Ulas', sname='Kamaci']{Ulas Kamaci}
\affiliation{ University of Illinois Urbana-Champaign \\
Department of Electrical and Computer Engineering, and Coordinated Science Laboratory \\
Urbana, IL 61801, USA
}
\email{ukamaci2@illinois.edu}

\author{Farzad Kamalabadi}
\affiliation{ University of Illinois Urbana-Champaign \\
Department of Electrical and Computer Engineering, and Coordinated Science Laboratory \\
Urbana, IL 61801, USA
}
\email{farzadk@illinois.edu}


% \author[orcid=0000-0000-0000-0001,sname='North America']{Tundra North America}
% \altaffiliation{Kitt Peak National Observatory}
% \affiliation{University of Saskatchewan}
% \email[show]{fakeemail1@google.com}  

% \author[orcid=0000-0000-0000-0002,gname=Bosque, sname='Sur America']{Forrest Sur Am\'{e}rica} 
% \altaffiliation{Las Campanas Observatory}
% \affiliation{Universidad de Chile, Department of Astronomy}
% \email{fakeemail2@google.com}

% \author[gname=Savannah,sname=Africa]{S. Africa}
% \affiliation{South African Astronomical Observatory}
% \affiliation{University of Cape Town, Department of Astronomy}
% \email{fakeemail3@google.com}

% \author{River Europe}
% \affiliation{University of Heidelberg}
% \email{fakeemail4@google.com}

% \author[0000-0000-0000-0003,sname=Asia,gname=Mountain]{Asia Mountain}
% \altaffiliation{Astrosat Post-Doctoral Fellow}
% \affiliation{Tata Institute of Fundamental Research, Department of Astronomy}
% \email{fakeemail5@google.com}

% \author[0000-0000-0000-0004]{Coral Australia}
% \affiliation{James Cook University, Department of Physics}
% \email{fakeemail6@google.com}

% \author[gname=IceSheet]{Penguin Antarctica}
% \affiliation{Amundsen–Scott South Pole Station}
% \email{fakeemail7@google.com}

% \collaboration{all}{The Terra Mater collaboration}

%% Use the \collaboration command to identify collaborations. This command
%% takes an optional argument that is either a number or the word "all"
%% which tells the compiler how many of the authors above the command to
%% show. For example "\collaboration[all]{(DELVE Collaboration)}" wil include
%% all the authors above this command.
%%
%% Mark off the abstract in the ``abstract'' environment. 
\begin{abstract}
\noindent
Although solar imaging spectroscopy has been widely used to study the Sun, the
inherent rastering process in conventional slit spectrographs required to generate 
spatial-spectral images limits their capability for studying dynamic events. In 
contrast, slitless imaging spectrographs can enable cotemporal imaging
and spectroscopy by removing the slit and the rastering process, at the expense 
of producing overlapped spectra on the observations, implying the necessity for
the solution of the resulting inverse problem. In this work, we develop a deep
learning approach to disentangle the overlapped spectra for a slitless
spectrograph design with a narrow bandpass in extreme ultraviolet (EUV) and
detectors placed at multiple orders. The bandpass of the instrument is assumed
to contain a single dominant spectral line, and our method recovers the three
Gaussian line parameters (intensity, Doppler velocity, and line width) from the
multi-order observations. We train a convolutional neural network (CNN) as a
mapping from the spectrograph observations to the spectral line parameters in a
supervised manner, using a large dataset derived from Hinode/EIS observations.
We compare the proposed method to a widely used tomographic reconstruction
approach, MART, and a recently proposed parametric inversion approach. Experiments on various observational settings show that the proposed
method substantially outperforms both baselines in intensity and velocity
reconstruction (velocity errors reduced by more than a factor of two relative to
the closest baseline), and is the only method to achieve
meaningful line-width recovery. Reconstruction is completed in under a second,
orders of magnitude faster than the alternative methods.

\end{abstract}

%% Keywords should appear after the \end{abstract} command. 
%% The AAS Journals now uses Unified Astronomy Thesaurus (UAT) concepts:
%% https://astrothesaurus.org
%% You will be asked to selected these concepts during the submission process
%% but this old "keyword" functionality is maintained in case authors want
%% to include these concepts in their preprints.
%%
%% You can use the \uat command to link your UAT concepts back its source.
\keywords{\uat{Ultraviolet spectroscopy}{2284} --- \uat{Solar corona}{1483} ---
\uat{Convolutional neural networks}{1938} --- \uat{Nonlinear regression}{1948}}

% \keywords{methods: statistical -- Sun: corona -- Sun: UV radiation --
% techniques: image processing -- techniques: neural networks --techniques:
% imaging spectroscopy}

%% From the front matter, we move on to the body of the paper.
%% Sections are demarcated by \section and \subsection, respectively.
%% Observe the use of the LaTeX \label
%% command after the \subsection to give a symbolic KEY to the
%% subsection for cross-referencing in a \ref command.
%% You can use LaTeX's \ref and \label commands to keep track of
%% cross-references to sections, equations, tables, and figures.
%% That way, if you change the order of any elements, LaTeX will
%% automatically renumber them.

\section{Introduction} \label{sec:intro}

Solar imaging spectroscopy is a fundamental technique to study the solar
corona, which provides spatial and spectral information of a two-dimensional 
field of view. It allows for inferring key physical parameters of the plasma 
such as the velocity, temperature, and density. The study of these parameters
sheds light on the complex dynamic behavior of solar plasma
\citep{phillips2012ultraviolet, rybicki2008radiative}. 


% \textit{Examples, applications (what is it used for - Doppler 
% shift, temperature, plasma density \dots)}

% Scientific Goal:
% \begin{itemize}
%   \item ESIS is designed to reveal the structure and dynamics of small-scale transient
%   events that are prevalent at transition region temperatures
% \end{itemize}

% Due to the inherent limitation of two-dimensional detectors to capture the 
% three-dimensional spatial-spectral data cube, traditional imaging spectrographs
% require a scanning process where a two-dimensional slice of the data cube is
% captured at each exposure. While filtergraph type spectrographs scan the cube in
% the wavelength dimension by swapping bandpass filters, dispersing slit
% spectrographs disperse light coming from a thin stripe of the field of view
% through the slit, which is then scanned spatially to build up the cube.
% % \textit{Define filtergraph-type vs dispersing slit spectrograph
% % with 1 sentence each.}

Traditional imaging spectrographs rely on a scanning process to construct the
three-dimensional spatial-spectral data cube. Two main classes of such
spectrographs are filtergraph type and dispersing slit type, where the former
scans the cube in the wavelength dimension by swapping bandpass filters at each
exposure, and the latter by dispersing light coming from a thin stripe of the
field of view through the slit, which is then scanned spatially to construct the
cube.

The main drawback of the filtergraph type spectrographs is the low spectral 
resolution they can provide that limits their ability to infer the density,
temperature and velocity accurately. Dispersing slit spectrographs on the other
hand suffer from large scanning time ($\sim 10$ minutes) that renders them
unsuitable for imaging dynamic events such as solar flares, coronal mass
ejections, and transient brightening that evolve on the orders of seconds
\citep{kankelborg2001simultaneous}.

% \item Scanning based spectrographs are not suitable for imaging dynamic events
% \textit{(give numbers, typical scanning time vs solar feature evolution time)(citation).}

\begin{figure*}[t]
\centering
\includegraphics[width=0.47\textwidth]{slitli.png} 
\hspace{0.05 in} 
% \vline 
\begin{tikzpicture}
    \draw[dashed] (0,0) -- (0,3.5);
\end{tikzpicture}
\hspace{0.00 in}
\includegraphics[width=0.47\textwidth]{slitless.png}
\caption{Illustration of measurements in slit and slitless spectrographs. The
left panel shows a slit spectrograph, which requires scanning the slit across the
field of view (FOV). The right panel depicts a slitless spectrograph, where
multiple detectors are positioned at diffraction orders -1, 0, 1 to increase
measurement diversity. In the slitless configuration, the absence of a slit
causes spectral overlap along the spatial dimension parallel to the dispersion
direction.}
\label{fig:slitless}
\end{figure*}

Slitless spectrographs are modifications of slit spectrographs where the slit is
completely removed to capture \textit{cotemporal} spatial and spectral
information of a scene with a large field of view, making them suitable for
studying dynamic events. However, their measurements are difficult to interpret
and analyze for extended sources like the Sun due to the overlapping spectrum of
the neighboring pixels on the detector in the dispersion direction (see Figure
\ref{fig:slitless} for illustration). Therefore, their measurements, also called
\textit{overlapograms}, need to be inverted to reveal the spectrum at each
spatial position in the field of view of the instrument. 

Slitless spectrographs can take measurements at multiple diffraction orders, 
where each order mathematically corresponds to a projection of the
three-dimensional spatial-spectral data cube $r(x,y,\lambda)$ on the
two-dimensional detector.  Therefore, they fall in the category of computed
tomography imaging spectrographs (CTIS) \citep{okamoto1991simultaneous,
descour1995computed} which are non-scanning spectrographs that take
multiple simultaneous projections of the data cube. CTIS systems are inherently
limited-angle tomography systems because they rely on a limited set of
spectrally dispersed projections, determined by the disperser design and the
size of the focal plane array, which affect the accuracy of the spectral
reconstruction.

The number of projections in CTIS systems plays an important role in the
fidelity of reconstructions. Early slitless spectrographs such as the S082A
instrument on Skylab \citep{tousey1977extreme} and the Res-K instrument of the
KRONOS-I mission \citep{zhitnik1998spectral} routinely captured only a single
projection of the data cube that made it impossible to invert their data and
limited the analysis. Two projections already enable the estimation of Doppler
shifts by analyzing the difference image between them \citep{mariska1992solar,
courrier2018using}. The Multi-Order Solar Extreme ultraviolet Spectrograph
(MOSES) \citep{fox2011snapshot} captures three projections at the diffraction
orders -1,0,1 that enables full inversions to reconstruct the data cube
\citep{fox2003data}. The more recent EUV Snapshot Imaging Spectrograph (ESIS)
\citep{parker2022first} captures four unique projections of the data cube for
more accurate estimations of the spectral line parameters.

% The S082A
% instrument on Skylab \cite{tousey1977extreme} and the Res-K instrument of the
% KRONOS-I mission \cite{zhitnik1998spectral} were the only two satellite missions
% that routinely captured solar slitless spectrograph data. Since then, there has
% been more slitless spectrograph missions being proposed and developed including
% the Multi-Order Solar Extreme ultraviolet Spectrograph (MOSES)
% (\cite{fox2011snapshot}) COronal Spectroscopic Imager in the EUV (COSIE)
% (\cite{winebarger2019unfolding}) and the EUV Snapshot Imaging Spectrograph
% (ESIS) (\cite{parker2022first}) as well as the multi-slit approach the Multi
% Slit Solar Explorer (MUSE) (\cite{de2019multi}). 

% Slitless spectrograph measurements mathematically correspond to a tomographic
% projection of the three-dimensional spatial-spectral data cube $r(x,y,\lambda)$
% on the two-dimensional detector. Early slitless spectrographs such as the S082A
% observed only one projection, which limited their ability to estimate the
% spectral line parameters except for very compact sources. Slitless spectrographs
% with two projections are able to estimate the Doppler shifts
% (\cite{courrier2018using}), and invert temperature and density information
% (\cite{winebarger2019unfolding}). As with any tomographic inversion problem,
% adding more projections significantly increases the fidelity of the
% reconstructions. The MOSES instrument
% (\cite{kankelborg2001simultaneous,fox2003data}) operates in a narrowband and
% employs three detectors at the diffraction orders -1,0,+1 to enable full
% inversions to disentangle the overlapped spectra in the measurements. The recent 
% ESIS instrument is a more advanced design than MOSES with four unique
% projections.

% A nonscanning spectrograph that takes multiple simultaneous projections of the
% data cube is called a computed tomography imaging spectrograph (CTIS)
% (\cite{descour1995computed}). Because of the limited size of the focal plane
% array, CTIS instruments are limited-angle tomographic systems. The tomographic 
% reconstruction problem they need to solve is ill-posed due to the limited view 
% angles they can collect which leads to an under-determined problem with a 
% non-trivial null space referred to as the \textit{missing cone}
% (\cite{descour1995computed}). As the under-determined problems have infinitely 
% many solutions, incorporating prior knowledge of the object plays a critical 
% role in constraining the solution space to physically feasible ones.

% Slitless imaging spectrographs have been widely used in astronomical
% imaging to analyze compact sources such as galaxies and stars and are still
% being used today on the Hubble and the James Webb space telescopes
% (\cite{ruiz2005probing, willott2022near}). However, when it comes to extended
% sources, interpretation and analysis of the slitless spectrograph data becomes
% challenging because of the overlapping spectra in the measurements.

Due to the limitations in the detector area and the feasible number of
projections that can be captured simultaneously, the tomographic reconstruction
problem that the solar slitless spectrographs need to solve is an
under-determined and ill-posed inverse problem. Therefore, the prior knowledge
and the constraints that are incorporated in the reconstruction algorithm play a
crucial role in the fidelity of reconstructions.
Various inversion methods have been applied to the CTIS systems operating at
different wavelength regimes. For solar applications, most algorithmic 
contributions came from inverting the MOSES observations. \citet{fox2003data} have
compared the performance of Fourier back-projection, multiplicative algebraic
reconstruction technique (MART), and the Pixon technique for inversion. Although
MART has been reported to perform the best among those methods, all of them
suffer from over-smoothing of the line profiles in the reconstructions, and
underestimation of the velocities \citep{fox2011snapshot}. Furthermore, the
MART reconstruction results presented in \citet{fox2003data} and
\citet{fox2011snapshot} for simulated MOSES data are optimistic since they assume
an additional projection of the actual cube is available as a prior that is in
reality unknown.

MART was first implemented on a CTIS system by \citet{okamoto1991simultaneous}
and is still the preferred inversion technique for the recent ESIS mission
\citep{parker2022first}. It is favored for being a simple and a fast 
technique that can also incorporate non-negativity constraints and smoothness
priors. However, MART has two significant drawbacks: first, it struggles to accurately
recover spectral line parameters, particularly velocity and line width; second,
it is difficult to integrate more powerful priors to improve the solutions.

To address the limitations of MART, \citet{oktem2014parametric} and
\citet{davila2019slitless} proposed a parametric inversion approach where they 
exploit the fact that the EUV spectra of emission lines can be parameterized by
Gaussians \citep{rybicki2008radiative, ireland2005precision}. This
representation automatically constrains the shape of the reconstructed profiles
to Gaussians, dramatically decreasing the number of unknowns in the data cube.
Although the mentioned approach is promising, their proposed iterative
inversion technique has two drawbacks. First, they set up the three-dimensional
tomographic inversion as a set of two-dimensional problems that are solved
independently, which restricts them to use only one dimensional spatial priors.
Second, they only use simple smoothness penalties as spatial priors in their
inversion. Therefore, there is need for a more powerful parametric inversion
technique that can incorporate better suited priors for solar scenes
\citep{davila2019slitless}.

Choosing the best suited prior or regularization is a fundamental pursuit in
ill-posed computational imaging problems literature. Initially, simple priors like
smoothness or sparsity were employed to constrain solutions, relying heavily on
assumptions about the underlying structure of the images. The advent of machine
learning, particularly deep learning, has revolutionized this field by enabling
the use of data-driven priors learned from large datasets. These approaches can
capture complex patterns and structures in images, leading to substantial
improvements in the accuracy and robustness of solutions to inverse problems
\citep{goodfellow2016deep, jin2017deep, ongie2020deep}.


% While enforcing smoothness prevents overfitting to noise, it
% oversmoothes the reconstructions and looses the details in the image. With the 
% advent of compressed sensing, sparsity based priors have been proven to be
% useful in ill-posed inverse problems. More recently, deep learning techniques 
% surpassed the performance of sparsity based methods thanks to their ability to 
% extract relevant features and learn implicit priors effectively from training 
% data.

% A deep neural network (DNN) is a universal function approximator
% (\cite{hornik1989multilayer}) that takes an input and applies a set of linear 
% and nonlinear transformations through its multiple layers between its input and
% output. They can model complex nonlinear relations between the input and output
% space. There are various ways to utilize DNNs to solve inverse problems
% \cite{ongie2020deep}. If the imaging forward model $\bm {\mathcal A}$ is known,
% i.e. $\bm y = \bm {\mathcal A} (\bm x)$, a straightforward way to utilize a DNN
% is to use it as a reconstruction network $f_{\theta}(\cdot)$ that takes as
% input the measurement $\bm y$ and outputs the reconstruction $\hat {\bm x}$ such
% that $\hat{\bm x} = f_{\theta}(\bm y)$. The training can be done in a supervised 
% manner with a simulated dataset of matched input output pairs $(\bm x, \bm y)$.

% \begin{itemize}
% \item \textit{The use of deep learning in similar imaging problems (possibly
% nonlinear). Cite some BASP work on astronomy (e.g. gravitational lensing,
% black hole imaging).}
% \end{itemize}
% Deep learning 

In this paper we combine the parametric inversion idea proposed by
\citet{davila2019slitless} with a deep learning approach to solve the slitless
reconstruction problem. The approach is based on supervised training of a
convolutional neural network (CNN) with a large dataset derived from Hinode/EIS
measurements. The network takes as input the spectrograph measurements, and
returns the spectral line parameters for every pixel in the field of view.
% There are several advantages of this approach over the previous approaches.
% Supervised training allows the CNN to learn the nonlinear inverse mapping from
% the spectrograph measurements to the Gaussian parameters. The learned mapping is
% inherently regularized by the training dataset which consists of actual spectral
% measurements. Hence, the network is able to capture the statistical and
% structural information in the training set to produce feasible solutions from
% limited measurements. 
Reconstruction results demonstrate the superiority of this learning based
approach compared to the classical tomographic reconstruction approach MART \citep[as
implemented by][]{parker2022first} and the parametric inversion approach by
\citet{davila2019slitless}, including meaningful line-width recovery that the
existing methods fail to achieve. Typical reconstruction time is also less than a
second once training is done, which is orders of magnitude less than the
alternative reconstruction methods.

The proposed method demonstrates that data-driven priors learned from real solar
observations substantially outperform classical regularization across all three
spectral line parameters. The improvement is especially pronounced for
line-width recovery: the classical baselines fail to track spatial variation in
line width across the field of view, while the proposed method recovers it
reliably. These results establish deep learning as a viable and superior
framework for data-driven slitless EUV spectral inversion.

The organization of the paper is as follows. Section \ref{sec:par_fwd_mod} describes the parametric
forward model characterizing the connection between the measurements and the data
cube. Section \ref{sec:learned_recon} describes the proposed learning based framework. Section \ref{sec:num_res}
presents experiments and numerical results. Section \ref{sec:disc} discusses the
findings and concludes.


% \section{Slitless Imaging Spectroscopy} \label{sec:slitless_spec}

% % By the end of this section, you will learn how the slitless spectrograph 
% % measurements are related to the input spectral cube. You will appreciate the 
% % ill-posedness of the inverse problem, and understand the motivation behind the 
% % parametric modeling.

% \begin{itemize}
% \item First example, Russian Skylab, single dispersion - full FOV, hard to analyze
% \item MOSES - First multi-order slitless design to enable spectral deconvolution
% \item Explosion of recent missions: EUNIS, MUSE, COSIE, SERTS, ESIS
% \item CTIS systems in visible spectrum for non-space applications
% \end{itemize}


% \begin{figure}[h!]
% \centering
% \includegraphics[width=0.45\textwidth]{moses_ab.png}
% \caption{Illustration of a slitless imaging spectrograph with -1, 0, 1 
% orders producing measurements of a bichromatic source.}
% \label{fig:moses_ab}
% \end{figure}

\section{Forward Model} \label{sec:par_fwd_mod}
We consider a multi-order slitless spectrograph model similar to MOSES
\citep{kankelborg2001simultaneous} where multiple detectors capture independent
projections at different diffraction orders simultaneously. Let $a \in
\{0,-1,1,-2,2,\dots\}$ denote the diffraction order. We consider $K$ detectors
placed at the distinct diffraction orders $a_1, a_2, \dots, a_K$.  Let
$r(x,y,\lambda)$ represent the diffraction-limited spatial-spectral data cube,
and $y^{(a)} (x',y')$ the projection of the data cube incident on the detector
at the diffraction order $a$.  For convenience, all the spatial coordinates
$(x,y,x',y')$ and the spectral coordinate $\lambda$ are expressed at the
detector plane in units of the detector pixel size.  The relation between the
data cube and the detector measurements can then be expressed as line integrals
similar to parallel beam tomography of an optically thin object.  Assuming that
the dispersion direction is parallel to the $y$ axis, the projection
angle $\theta$ lies on the $(y,\lambda)$ plane and is determined by the
diffraction order $a$ according to $\theta = \tan^{-1}(a)$.  The governing
integral equation for our slitless spectrograph can be written as
\citep{kankelborg2001simultaneous}:
\begin{align}
    % y^{(a)}_{ij} & = \sum_k r_{i(j+ak)k} \\
    y^{(a)}(x',y') & = \int_{\lambda_0-\delta}^{\lambda_0+\delta} r(x', y' -
    a\lambda+a\lambda_0,\lambda) \ d\lambda 
\label{eq:fwd_cont}
\end{align}
where the integration paths are parameterized by $\lambda$ with a slope of $-a$ 
in the $(y,\lambda)$ plane, and are contained within the bandpass of the
instrument that has a width of $2 \delta$ centered around the rest wavelength
$\lambda_0$. Note that the linear model in Eq. \eqref{eq:fwd_cont} 
assumes that the detector is not saturated.

To express the discrete version of the model in Eq. \eqref{eq:fwd_cont}, let us
define the discrete versions of the data cube $r_{ijk} \in \mathbb R^{M \times M
\times P}$ and the measurements $y^{(a)}_{ij} \in \mathbb R^{M \times M}$
where $(i,j) \in \{1,\dots,M\}$ are spatial indices and $k \in \{1,\dots,P\}$ is
the spectral index. Each detector has $M \times M$ pixels, and the data cube is 
discretized with the same spatial resolution as the detector, leading to the
same size of $M \times M$ in the spatial coordinates. Typically, $P \ll M$ due 
to the narrow bandpass of the instrument. With this, the discrete forward model 
can be written as the weighted sum of the data cube voxels:
\begin{align}
    y^{(a)}_{i'j'} & = \sum_{j} \sum_{k} \beta^{(a)}(j, k;i', j') r_{i'jk}
\label{eq:fwd_sum}
\end{align}
where $\beta^{(a)}(j, k;i', j')$ represents the weight for the contribution of the
voxel at the position $(j,k)$ to the detector pixel at $(i',j')$ for the
diffraction order $a$.
We can express this linear model as a single matrix vector product by defining 
the measurement vector $\bm y \in \mathbb R^{KM^2}$ consisting of all detector 
pixels concatenated, and the data cube vector $\bm r \in \mathbb R^{PM^2}$ as:
\begin{equation}
    \bm y = \bm H \bm r
\end{equation}
where $\bm H \in \mathbb R^{KM^2 \times PM^2}$ is the forward matrix 
which is the undersampled version of the Radon transform matrix due to the
limited angle coverage of the projections.

To account for the detector noise, we add a Poisson noise generator process
$\mathcal P_\gamma(\cdot)$ to the forward model with a signal-to-noise
ratio (SNR) of $\gamma = \sqrt{N_{\text{avg}}}$ where $N_{\text{avg}}$
corresponds to the average number of photon counts per pixel for each detector.

% which is defined as the SNR corresponding to the average 
% number of counts ($\gamma = \sqrt{N_{\text{avg}}}$) per pixel per detector.

% , which is affected by many factors such as the scene
% brightness, instrument sensitivity, and exposure time. Since SNR changes with 
% photon count, we define the SNR per each detector 
% SNR also changes with
% position on the detector since the scene brightness changes. For each detector,
% we define $\gamma$ to be the SNR value corresponding to the pixel with the $90\%$
% of the maximum pixel value in that detector, providing a representative estimate
% of high signal regions while excluding extreme outliers. 

The forward model with the noise included becomes: 
\begin{equation}
    \bm y = \mathcal{P}_\gamma(\bm H \bm r)
\end{equation}

\subsection{Parametric Spectrum}
We model the instrument's passband to be dominated by a single bright EUV line
of interest, alongside secondary spectral blends and a background continuum.
This emission line is parameterized as a Gaussian profile, which conveniently
captures the three plasma parameters of interest: intensity, Doppler shift, and
the line width. These parameters vary with spatial position $(i,j)$ and are
denoted by $(f_{ij}, \mu_{ij}, \sigma_{ij})$, respectively, where $\mu_{ij}$ and
$\sigma_{ij}$ are expressed in units of detector pixels along the dispersion
direction. With that, each voxel of the data cube $r_{ijk}$ can be represented
by:
\begin{equation}
r_{ijk} 
= \frac{f_{ij}}{\sigma_{ij}\sqrt{2\pi}}
e^{-\frac{(k-k_0-\mu_{ij})^2}{2\sigma_{ij}^2}} + s_{ijk} + b_{ij}
\label{eq:param_voxel}
\end{equation}
where $k_0$ is the location of the rest wavelength of the emission line on the
discretized wavelength grid, the term $s_{ijk}$ contains any contribution from
secondary contaminant lines within the passband, and the term $b_{ij}$ is a
spectrally flat background amplitude that is constant across wavelength bins.

While this parametric formulation significantly reduces the dimensionality of
the inverse problem, it remains consistent with the physical characteristics of
typical EUV observations. The assumption of a single dominant line is
validated by the differential emission measure (DEM) analysis for the MUSE
mission \citep{de2019multi}, which demonstrates that a proper passband selection
can isolate a dominant EUV line such that its intensity exceeds secondary
contaminant lines by an order of magnitude or more.  Furthermore, although solar
plasmas in regions such as flares can exhibit physical deviations from a purely
Gaussian profile, such as heavier tails \citep{lorinvcik2020plasma}, current EUV
spectrographs possess instrumental widths comparable to the inherent widths of
coronal lines; the resulting convolution with the instrument response function
inherently smooths these physical deviations, rendering the observed profiles
highly Gaussian. 

\subsection{Parametric Forward Model}
To formulate the forward model parametrically across the entire field of view,
let $\bm x = [\bm f^\top \ \bm \mu ^\top \ \bm \sigma^\top ]^\top \in \mathbb
R^{3M^2}$ denote the vectorized Gaussian parameters of our emission line, $\bm z
= [\bm z_1^\top \ \bm z_2 ^\top \cdots \bm z_C^\top ]^\top \in \mathbb
R^{3CM^2}$ the Gaussian parameters of the secondary line blends where $\bm z_c =
[\bm f_c^\top \ \bm \mu_c ^\top \ \bm \sigma_c^\top ]^\top \in \mathbb
R^{3M^2}$, and $\bm b \in \mathbb R^{M^2}$ the additive constant background term
for each point in the $M \times M$ FOV. The data cube $\bm r$ can then be
expressed using the Gaussian data cube generator operators $\bm \Psi_c: \mathbb
R^{3M^2} \mapsto \mathbb R^{PM^2}$ for $c = 0, 1, \ldots, C$ as:
\begin{equation}
\bm r = \bm \Psi_0 (\bm x) + \sum_{c=1}^{C} \bm \Psi_c (\bm z_c) + \bm b
= \bm \Psi(\bm x, \bm z, \bm b)
\end{equation}
where $\bm b \in \mathbb R^{M^2}$ is broadcast across the $P$ spectral bins when
added to $\bm r \in \mathbb R^{PM^2}$.
Applying the projection operator $\bm H$ to this parametric form, the forward
model is expressed as:
\begin{equation}
\bm y = \mathcal{P}_\gamma(\bm H \bm r) = \mathcal{P}_\gamma(\bm H \bm \Psi (\bm
x, \bm z, \bm b))
\end{equation}
which can be written in a more compact form by defining the composite operator
$\bm{\mathcal A} := \bm H \circ \bm \Psi$:
\begin{equation}
\bm y = \mathcal{P}_\gamma(\bm {\mathcal A}(\bm x, \bm z, \bm b))
\label{eq:fwd_nl}
\end{equation}

To highlight the differences between the discretized and the parametric forms,
note that while $\bm H$ is linear, $\bm {\mathcal A}$ is a nonlinear forward
operator. Also note that while the parameterized form $(\bm x, \bm z, \bm b)$
has $(3C+4)M^2$ parameters, the full cube $\bm r$ has $PM^2$ parameters, where
typically $P \gg 3C+4$. Therefore, the parametric model results in fewer unknowns
in the data cube, at the cost of leading to a nonlinear forward model. The
model can be extended to include additional spectral components or non-Gaussian
profiles, though increased spectral complexity reduces identifiability given the
limited number of projections $K$.

\subsection{Inverse Problem}

The goal of the inverse problem is to retrieve the target spectral line
parameters $\bm x$ from the detector measurements $\bm y$. Because the
projection operator $\bm H$ captures only a limited number of view angles, the
mapping from scene to measurements is underdetermined, rendering the inverse
problem ill-conditioned and the solution susceptible to noise amplification.

A traditional approach is to first recover the full data cube $\bm r$ directly
from the measurements by solving a volumetric regularized least-squares problem:
\begin{equation}
\hat{\bm r} = \argmin_{\bm r} \|\bm y - \bm H \bm r\|_2^2 + \mathcal R(\bm r)
\label{eq:rls_vol}
\end{equation}
Because $\bm H$ is linear, Eq.~\eqref{eq:rls_vol} is a linear inverse problem.
Sparsity-promoting choices of $\mathcal R(\bm r)$, such as total variation or
$\ell_1$-based penalties, can be minimized directly via convex optimization
\citep{sidky2008image, rudin1992nonlinear} and are widely used in limited-angle
tomographic reconstruction. Algebraic methods such as ART and SART
\citep{gordon1970algebraic, andersen1984simultaneous} take a different approach:
rather than minimizing Eq.~\eqref{eq:rls_vol} to convergence, they perform
successive projection-based corrections that approximately satisfy the
measurement constraints, making them computationally efficient and well-suited
to the limited-angle setting. MART \citep{gordon1970algebraic} similarly seeks
measurement consistency through multiplicative updates, but is associated with
KL-divergence minimization rather than least-squares \citep{byrne1993iterative}.
However, recovering $\bm r \in \mathbb R^{PM^2}$ from $\bm y \in \mathbb
R^{KM^2}$ with $K \ll P$ remains severely underdetermined regardless of solver,
and hand-crafted regularizers $\mathcal R(\bm r)$ struggle to capture the
complex morphology of coronal spectral cubes. Moreover, since the ultimate goal
is to retrieve the primary emission line parameters $\bm x$ rather than the full
spectral cube $\bm r$, an additional spectral decomposition step is needed to
extract $\bm x$ from the recovered $\hat{\bm r}$:
\begin{equation}
\hat{\bm x}, \hat{\bm z}, \hat{\bm b} = \argmin_{\bm x, \bm z, \bm b}
\|\bm \Psi(\bm x, \bm z, \bm b) - \hat{\bm r}\|_2^2
\label{eq:spec_fit}
\end{equation}
This sequential pipeline, volumetric inversion followed by spectral fitting,
is not jointly optimized, introducing a compounded source of estimation error.

The parametric approach avoids this sequential pipeline by directly optimizing
for the spectral parameters. Prior formulations
\citep{oktem2014parametric, davila2019slitless} assumed a single Gaussian
component and neglected secondary blends and background; here we adopt the full
multi-component model of Section \ref{sec:par_fwd_mod}, restricting $\bm r$ to
the space of physically realistic spectral profiles:
\begin{equation}
\hat{\bm x}, \hat{\bm z}, \hat{\bm b} = \argmin_{\bm x, \bm z, \bm b}
\|\bm{\mathcal A}(\bm x, \bm z, \bm b) - \bm y\|_2^2 +
\mathcal R(\bm x, \bm z, \bm b)
\label{eq:rls_par}
\end{equation}
The reduced parameter space of $(3C+4)M^2 \ll PM^2$ unknowns substantially
improves identifiability, and the parametric formulation implicitly encodes a
strong physics-based prior on spectral shape: by construction, only solutions
consistent with Gaussian-profile spectral lines are admitted, a constraint that
hand-crafted volumetric regularizers cannot readily replicate. The cost is a
nonlinear forward operator $\bm{\mathcal A}$, which renders the optimization
non-convex; convergence to the global minimum is not guaranteed, and the
solution is sensitive to initialization and the choice of regularizer.

Both classical approaches require solving a high-dimensional optimization
problem for each new observation, which is computationally expensive and limits
their applicability to large-scale data. Furthermore, designing effective
regularizers $\mathcal R$ does not readily adapt to the spatial diversity of
observed coronal scenes. These limitations motivate a data-driven approach,
described in Section \ref{sec:learned_recon}, which replaces the per-observation
optimization with a learned mapping trained once on a large corpus of realistic
observations.

\section{Learned Reconstruction} \label{sec:learned_recon}
In this section we describe the developed learning based approach to
reconstructing the data cube from the slitless spectrograph measurements.

\subsection{Learning Approach}
In contrast to the classical approaches of Section \ref{sec:par_fwd_mod}, which
require solving a high-dimensional optimization problem for each new observation,
we adopt a learning-based approach that replaces per-observation inversion with a
direct data-driven mapping from measurements to spectral parameters. A deep
reconstruction network $f_\theta: \mathbb R^{KM^2} \rightarrow \mathbb R^{3M^2}$
is trained to directly map the dispersed measurements $\bm y$ to the target
spectral parameters $\bm x$, bypassing both the intermediate volumetric
reconstruction step and the subsequent spectral fitting stage required by
volumetric methods. The network weights are optimized once over a large dataset
$\mathcal D$ of matched $(\bm x, \bm y)$ pairs:
\begin{equation}
\hat {\theta} = \argmin_{\theta} \sum_{(\bm x_i, \bm y_i) \in \mathcal D}
\mathcal{L}(\bm x_i, f_\theta (\bm y_i))
\label{eq:erm}
\end{equation}
Here $\mathcal{L}$ is a task-specific loss function detailed in Section
\ref{sec:training}. Once trained, $f_\theta$ is fixed and applies to new
observations via a single forward pass, making inference orders of magnitude faster than classical
per-observation solvers. Rather than relying on hand-crafted regularizers, the
network implicitly learns spatial and spectral priors for each parameter from
data, capturing the complex morphology of coronal structures more effectively
than analytic penalties. Furthermore, because the training measurements
$\bm y$ are derived directly from observational data that natively contains
secondary spectral blends and coronal background, the network learns to isolate
the primary line parameters $\bm x$ without requiring these components to be
explicitly modeled in the inversion.

\subsection{Dataset} \label{sec:dataset}
The adopted supervised training approach requires a dataset of matched pairs of
target parameters $\bm x$ and corresponding dispersed measurements $\bm y$. Our
approach for generating the dataset is to first obtain a collection of
spatial-spectral data cubes of the target spectral window from previous observations, which are then
interpolated onto our instrument's spectral grid to yield $\bm r$. We fit
Gaussians to these regridded cubes to obtain the ground-truth target parameters
$\bm x$, while simultaneously simulating the measurements $\bm y =
\mathcal{P}_\gamma(\bm H \bm r)$ by applying the linear projection operator to
the same cubes. Performing both the fitting and the projection on a common
observational cube ensures that $\bm x$ and $\bm y$ remain spectrally
consistent, and the resulting measurements natively incorporate the physics of
the EUV spectra, including faint secondary line blends, coronal background, and
non-Gaussian profile variations. This ensures the network learns to filter
these physical complexities robustly.

We utilize observations from the Extreme Ultraviolet Imaging Spectrometer (EIS)
aboard the Hinode observatory \citep{culhane2007euv}. Launched in September
2006, EIS observes the solar corona and upper transition region emission lines
in the wavelength ranges of \SIrange{170}{210}{\angstrom} and
\SIrange{250}{290}{\angstrom} with a spectral resolution of
\SI{22}{\milli\angstrom} and a plate scale of \SI{1}{\arcsecond} per pixel.
While the instrument offers four slit or slot widths, we restrict our analysis
to the highest spatial resolution scans.

To construct a comprehensive training dataset, we queried the Hinode Science
Data Centre (SDC) for EIS raster scans captured between December 2006 and
December 2025. To ensure sufficient spatial context and adequate sampling of our
target spectrum, the query was restricted to Active Region (AR) observations
utilizing the \SI{1}{\arcsecond} slit, requiring a field of view exceeding $64
\times 64$ arcseconds and capturing the Fe XII \SI{195.12}{\angstrom} spectral
window. This strict filtering yielded 2,383 candidate scans.

\begin{figure}[t]
\centering
\includegraphics[width=0.47\textwidth]{eis_fit.png}
\caption{Two-component spectral fit to a Hinode/EIS observation of the
Fe~\textsc{xii} \SI{195.12}{\angstrom} spectral window, used to extract
ground-truth Gaussian parameters for the training dataset.  \textit{Left:}
Summed intensity map of the spatial cutout; the black cross marks the pixel
shown in the right panel.  \textit{Right:} Observed intensity at the selected
pixel (black markers), decomposed into a combined model profile (blue), the
primary Fe~\textsc{xii} component (red solid), a secondary spectral blend (red
dashed), and a constant background (green). The primary Gaussian parameters,
intensity, Doppler velocity, and line width, are isolated from this fit to form
the clean ground-truth maps $\bm{x}$.}
\label{fig:eis_fit}
\end{figure}


The selected HDF5 files were processed using the EIS Python Analysis Code
(EISPAC) package \citep{weberg2023eispac}. We utilized the
\lstinline|eispac.read_cube| routine to apply basic pointing corrections and
pre-flight radiometric calibration, extracting the data cube centered around
the \SI{195.12}{\angstrom} wavelength. The corrected cube was then interpolated
onto our instrument's spectral grid, producing the common cube $\bm r$ from
which both the ground-truth parameters and the dispersed measurements are
derived. Because this spectral window contains a known secondary blend and
coronal background emission, the spectra were explicitly modeled using a template comprising two Gaussians
and a constant background (i.e., $C=1$ secondary component in the parametric
model of Section \ref{sec:par_fwd_mod}), as illustrated in Figure
\ref{fig:eis_fit}. This multi-component fitting was performed on $\bm r$
by solving Eq.~\eqref{eq:spec_fit} via the Levenberg-Marquardt algorithm \citep{levenberg1944method} due to its
flexibility and extensive legacy in solar and heliophysics research. The
parameters of the primary Fe XII Gaussian component were then isolated and
extracted to serve as the clean ground-truth two-dimensional maps of intensity,
velocity, and line width ($\bm x$). The same cube was concurrently fed through
the linear forward tomographic operator to simulate the dispersed measurements
($\bm y$) at five diffraction orders ($0, \pm1, \pm2$). Figure \ref{fig:dataset} contains example ($\bm x, \bm y$) pairs
from the dataset.

\begin{figure}[t]
\centering
\includegraphics[width=0.45\textwidth]{dset_compact.png}
\caption{Sample pairs from the EIS dataset for $K=3$ projections. Each column
represents a sample consisting of an $(\bm x, \bm y)$ pair: the first three rows
show the intensity, velocity, and line width of the parameters $\bm x$, while
the last three rows display the corresponding $0$, $-1$, and $1$ order
measurements of $\bm y$.}
\label{fig:dataset}
\end{figure}

To ensure the neural network is trained exclusively on physically valid and
high-fidelity signals, we applied a strict masking protocol to eliminate
anomalous data, cosmic ray hits, and non-convergent fits. Spatial locations were
masked and excluded from the dataset if the fit failed to converge, or if the
extracted parameters fell outside physically expected bounds: intensity outside
the range of \num{100} to \num{25000}
\unit{\erg\per\centi\metre\squared\per\second\per\steradian}, absolute velocity
exceeding \SI{68.5}{\kilo\metre\per\second}, or line width narrower than the
instrument's spectral resolution of \SI{0.022}{\angstrom}. To accommodate the
memory constraints of the neural network while simultaneously augmenting the
total number of training samples, we extracted $15$ independent $64 \times 64$
spatial patches from the valid, unmasked regions of each full field-of-view
observation.

The distributions of the ground-truth parameters across the resulting dataset
are shown in Figure \ref{fig:dset_hist}. Doppler velocities and line widths are
expressed in units of detector pixels along the dispersion axis; given the EIS
spectral dispersion of \SI{22}{\milli\angstrom\per\pixel}, one pixel
corresponds to approximately \SI{34}{\kilo\metre\per\second} of Doppler shift
at the Fe XII \SI{195.12}{\angstrom} rest wavelength. The intensity
distribution is strongly right-skewed with the bulk of pixels at lower values,
the velocity histogram is roughly symmetric about zero, as expected for an
ensemble of active-region pixels containing a mix of upflows and downflows, and
the line widths are clustered near \num{1.25} pixels.

To properly train and evaluate the network, the compiled dataset is partitioned
into training, validation, and test subsets, serving to optimize model learning,
tune hyperparameters, and provide unbiased evaluation, respectively. We employed
an \SI{80}{\percent}--\SI{10}{\percent}--\SI{10}{\percent} split for these
subsets. Crucially, to prevent data leakage, this partitioning was performed
at the level of unique parent observation rasters rather than individual
patches. This strategy ultimately yields \num{23036} patches for training
and \num{2880} for validation. The final \SI{10}{\percent} of the independent
parent rasters are held out as an unseen candidate testing pool, from which
representative patches are drawn exclusively for the comparative experiments
in Section~\ref{sec:num_res}.

\begin{figure}[t]
\centering
\includegraphics[width=0.47\textwidth]{dset_histograms.png}
\caption{Percentage-normalized histograms of the ground-truth Gaussian
parameters across the compiled dataset: intensity (left), Doppler velocity
(middle), and line width (right), aggregated over all training, validation, and
test patches. Velocity and line width are expressed in detector pixels along
the dispersion direction.}\label{fig:dset_hist}
\end{figure}

Prior to training the neural network, the ground-truth parameters ($\bm x$) and
the simulated measurements ($\bm y$) are normalized to ensure stable training
and balanced gradient contributions across channels. The intensity maps ($\bm f$) and the dispersed measurements ($\bm y$) are each
normalized by taking the natural logarithm of each value and then subtracting
the global mean and dividing by the global standard deviation, both computed
from the log-transformed training set. Both are normalized under the same type
of transform to maintain consistent input and output scales for the network,
though each uses its own independently computed statistics. The
log transform is motivated by the strongly right-skewed positive distribution of
coronal intensities (Figure \ref{fig:dset_hist}): by compressing the dynamic
range, it maps the intensity distribution to a near-symmetric form that is more
amenable to gradient-based optimization and equates the relative contribution of
errors at any intensity level. The velocity ($\bm \mu$) and line width ($\bm \sigma$) parameter maps
are normalized to zero mean and unit variance using the global training-set
statistics. Following the network's prediction, the inverse transformations are
applied to retrieve all parameters in their true physical units.
 
\begin{figure}[t]
\centering
\includegraphics[width=0.45\textwidth]{unet.png}
\caption{Diagram of the U-Net architecture used in this work for $K=3$ input
channels, corresponding to measurements from three detectors. A separate U-Net
is trained for each value of $K$, with the number of input channels matching the
number of detectors. In the encoder path, the number of channels doubles while
the spatial resolution halves after each layer, capturing increasingly abstract
features. In the decoder path, the number of channels halves while the spatial
resolution doubles, progressively reconstructing the output. Skip connections
between corresponding layers in the encoder and decoder paths preserve spatial
information.}
\label{fig:unet}
\end{figure}

% The signal-to-noise ratio (SNR) is set to \SI{35}{\decibel} in each detector
% channel $\bm y_i \in \mathbb R ^{64 \times 64}, \ i \in \{1,\dots,K\}$.

\subsection{Training and Network Architecture} \label{sec:training}
To learn the approximate inverse mapping $f_\theta$, we employ a Convolutional
Neural Network (CNN) based on the U-Net architecture \citep{ronneberger2015u}.
Originally developed for image segmentation, U-Net has become a foundational
architecture for complex image-to-image regression tasks \citep{isola2017image,
jin2017deep, ulyanov2018deep, zhou2019unet++}. Its efficacy stems from its
symmetric encoder-decoder structure, which extracts deep semantic features via
successive downsampling while recovering high-resolution spatial details through
upsampling and skip connections. This multi-scale feature representation
\citep{ye2019understanding, liu2022learning} is uniquely suited for our
tomographic inversion, as it allows the network to simultaneously analyze
localized spectral line profiles and broader regional coronal structures.

The network architecture is illustrated in Figure~\ref{fig:unet}. The network
takes a three-dimensional tensor as input, consisting of two spatial dimensions
and a channel dimension. We utilize the channel dimension to stack the $K$
dispersed spectral orders ($k \in \{1,\dots,K\}$) captured by the
simulated detectors ($\bm y$). The network outputs a tensor of identical spatial
dimensions, where the three channels correspond to the retrieved parameter maps:
intensity, velocity, and line width ($\bm{x}$). Within the network, the encoder
path reduces spatial dimensions using $2 \times 2$ max pooling operations, while
the decoder path recovers them via bilinear interpolation. Each downsampling and
upsampling step is separated by a double convolutional block consisting of two
sequential $3 \times 1$ convolutional layers, each followed by batch
normalization and a Rectified Linear Unit (ReLU) activation. The number of
feature channels progressively doubles at each encoder stage and halves at each
decoder stage.

For the loss function $\mathcal{L}$ in Eq.~\eqref{eq:erm} we adopt the mean squared
error (MSE). Because the log z-score normalization described above maps the
intensity distribution to a near-symmetric form with comparable variance to the
z-scored velocity and line width channels, plain MSE provides balanced gradient
contributions across all three output channels without requiring additional
per-channel reweighting.

The network weights are optimized by minimizing the empirical risk in
Eq.~\eqref{eq:erm} using the Adam optimizer \citep{kingma2014adam} with a learning
rate of $2 \times 10^{-4}$. We train the model for $200$ epochs with a batch
size of $32$. The architecture is implemented in the PyTorch framework, and the
training process completes in approximately $2.5$ hours on a single NVIDIA TITAN
Xp GPU.

\section{Numerical Results} \label{sec:num_res}
In this section, we evaluate the effectiveness of the proposed approach under
varying observational settings and compare it against two classical baselines:
MART, the established tomographic inversion method in the solar slitless
spectroscopy literature, and the 1D MAP approach of
\citet{davila2019slitless}, whose parametric formulation motivated the forward
model adopted in this work. The experiments vary the number of detector
projections $K$ and the signal-to-noise ratio $\gamma$, reflecting the key
practical design parameters of any slitless spectrograph: the number of
detectors determines $K$, and the exposure time and source brightness govern
$\gamma$. In all experiments, the instrument is modeled with a narrow bandpass
centered on the Fe XII \SI{195.12}{\angstrom} emission line, using the same
spectral dispersion scale of \SI{22}{\milli\angstrom\per\pixel} as the EIS
instrument.

For the projection count experiment, we consider the spectral orders
$\mathcal{S} = \{0, -1, +1, -2, +2\}$ and vary the number of detectors $K \in
\{2, 3, 4, 5\}$, placing them at the first $K$ orders in $\mathcal{S}$ (e.g.,
for $K = 2$, we use orders $0$ and $-1$). This evaluation is performed without
adding synthetic noise ($\gamma = \infty$), so the algorithms are tested on the
authentic Poisson noise already present in the raw EIS source data and
propagated through the forward model. For the noise experiment, we evaluate
performance across decreasing SNR levels ($\gamma \in \{30, 20, 10\}$) at
fixed $K = 3$, simulating lower-fidelity observational conditions such as
shorter exposure times or dimmer coronal regions.

While the proposed U-Net can evaluate thousands of patches in seconds, executing
1D MAP and MART across a large dataset is computationally intractable due to
their iterative bottlenecks. Consequently, all methods are evaluated on a
randomly selected subset of 100 independent $64 \times 64$ patches drawn from
the held-out testing pool, whose parameter distributions are representative of
the full EIS dataset. Measurements for each patch are simulated according to
the specific $(K, \gamma)$ configuration under evaluation.

The root mean square error (RMSE) and the systematic bias (mean of the signed
residuals) between the ground-truth parameters and the reconstructions serve as
the primary performance metrics. RMSE quantifies the overall magnitude of
reconstruction error, while bias reveals whether a method systematically over-
or under-estimates a parameter. All metrics are reported in physical units.

For a fair comparison, each method is separately optimized for each
observational setting $(K, \gamma)$. Because the number of input channels to
the U-Net equals $K$ and the noise statistics of the training measurements
depend on $\gamma$, a separate network is trained for each $(K, \gamma)$ pair
using the full EIS training dataset. The tunable hyperparameters of 1D MAP and
MART are likewise optimized for each setting, using the average reconstruction
RMSE over 50 randomly drawn patches from the EIS training set as the criterion.

\subsection{Compared Methods}
Here we describe the two classical baseline methods: the parametric
column-wise inversion approach of \citet{davila2019slitless}, hereafter referred
to as 1D MAP, and the tomographic inversion method MART.

\noindent
\textbf{1D MAP.} \ Following \citet{davila2019slitless} and
\citet{oktem2014parametric}, the 1D MAP approach formulates the inversion as a
regularized nonlinear least-squares optimization problem, equivalent to a
Maximum A Posteriori (MAP) estimation, that is solved independently for each
one-dimensional detector column:
\begin{equation}
\hat {\bm \xi}_i = \argmin_{\bm \xi_i} \|\bm y_i - \bm{\mathcal A}_i (\bm \xi_i) \|_2^2
+ \mathcal R (\bm \xi_i) \ , \quad i \in \{1,\dots,M\}
\label{eq:1dmap}
\end{equation}
where $\bm y_i \in \mathbb R^{KM}$ is the $i^{\text{th}}$ column of the
measurements and $\bm{\mathcal A}_i$ is the column-wise forward operator. The
original formulation of \citet{davila2019slitless} modeled a single Gaussian
component per pixel with no background term, giving $\bm \xi_i \in \mathbb
R^{3M}$ with three parameters per pixel: intensity, Doppler velocity, and line
width. Because our ground-truth parameters are derived by jointly
fitting the primary Fe XII component, a secondary spectral blend, and a
background term (Section \ref{sec:dataset}), a single-component forward model
would conflate background emission and the secondary blend with the primary
Gaussian, biasing the recovered intensity, velocity, and line width. We
therefore extend the 1D MAP to use the same two-component plus background
spectral model as the ground-truth fitting, providing a fair comparison. Each
column now optimizes a combined parameter vector
$\bm \xi_i = [\bm x_i^\top \ \bm z_{1,i}^\top \ \bm b_i^\top]^\top \in
\mathbb R^{7M}$, where $\bm x_i = [\bm f_i^\top \ \bm \mu_i^\top \ \bm
\sigma_i^\top]^\top \in \mathbb R^{3M}$ are the primary Gaussian parameters,
$\bm z_{1,i} = [\bm f_{1,i}^\top \ \bm \mu_{1,i}^\top \ \bm
\sigma_{1,i}^\top]^\top \in \mathbb R^{3M}$ are the secondary blend parameters,
and $\bm b_i \in \mathbb R^M$ is the per-pixel background amplitude; only $\bm
x_i$ is returned as the reconstruction output. The spatial smoothness prior is
applied to all seven parameter maps:
\begin{equation}
\mathcal R(\bm \xi_i) = \sum_{j=2}^{M} \sum_{p=1}^{7} w_p
\bigl(\xi_{i,j}^{(p)} - \xi_{i,j-1}^{(p)}\bigr)^2
\label{eq:1dmap_reg}
\end{equation}
where $p=1,2,3$ index the primary parameters $\bm x_i$, $p=4,5,6$ the
secondary blend parameters $\bm z_{1,i}$, and $p=7$ the background amplitude
$\bm b_i$. The per-parameter weights $w_p$ are tunable hyperparameters
optimized on a subset of 50 images of the EIS training set.

A fundamental weakness of this approach is its independent column-wise
formulation; reconstructing 1D columns via Eq. \eqref{eq:1dmap} ignores all
morphological coupling across the orthogonal spatial dimension. Additionally,
because the objective function is highly nonconvex, finding the global minimum
is computationally challenging. Although \citet{oktem2014parametric} developed a
dynamic programming algorithm to guarantee global optimality for the original
single-Gaussian, no-background 1D formulation, this guarantee does not extend
to the extended model used here, and the accuracy of even the original global
solution is inherently constrained by the simplistic 1D spatial prior.
Furthermore, its high computational cost makes it impractical for routine
processing of large-scale observational data.  Consequently, practical
implementations rely on local iterative solvers where convergence to a global
minimum is not guaranteed and depends heavily on initialization. While
\citet{davila2019slitless} originally implemented this using the Nelder-Mead
simplex algorithm \citep{nelder1965simplex}, we instead utilize the L-BFGS-B
method \citep{zhu1997algorithm} via the \lstinline|scipy.optimize.minimize|
Python library, as we found this quasi-Newton approach yields superior
convergence for this specific inverse problem.

\noindent
\textbf{MART.} \ The multiplicative algebraic reconstruction technique (MART)
has been developed for limited angle tomography problems
\citep{gordon1970algebraic, okamoto1991simultaneous} and has the
advantages of being a simple and fast method. It was reported to be the best
performing inversion algorithm among various others for the MOSES inversions
\citep{fox2011snapshot}. The recent ESIS mission also uses this method for
their inversions \citep{parker2022first}.

MART is based on a multiplicative update rule that aims to iteratively adjust
the reconstructed volume $\bm r \in \mathbb R^{PM^2}$ until its projections $\bm
H \bm r$ match the actual measured projection data $\bm y\in \mathbb R^{KM^2}$.
We implement an extended version of the MART variant described in
\cite{parker2022first}, adapted to use the two-component plus background
spectral model, as follows:

\begin{enumerate}
\item Construct the initial cube
\begin{align}
\bm r_0 =& \bm \Psi_0(f_1 \bm y^{(0)},\; \mu_1^0 \bm 1,\; \sigma_1^0 \bm 1)
\nonumber\\
         +& \bm \Psi_1(f_2 \bm y^{(0)},\; \mu_2^0 \bm 1,\; \sigma_2^0 \bm 1)+ f_{\rm bg}\, \bm y^{(0)}
\nonumber\\
\label{eq:mart_init}
\end{align}
where $f_1, f_2, f_{\rm bg}$ are the fractional amplitude contributions of the
primary component, secondary blend, and background (satisfying $f_1 + f_2 +
f_{\rm bg} = 1$); $\mu_c^0$ is the rest-wavelength centroid of component $c$
(zero Doppler offset); $\sigma_c^0$ is its nominal line width from the EIS
training set; and the last term is the scalar background map
$\bm b_0 = f_{\rm bg}\,\bm y^{(0)} \in \mathbb R^{M^2}$, broadcast across the
spectral dimension as in Section \ref{sec:par_fwd_mod}.
\item Enhance the contrast of $\bm r$ and normalize:
\begin{equation}
\bm r \leftarrow \frac{\bm r+\bm r^{(1+\alpha)}}{\sum_{xy\lambda}\bm r + \bm
r^{(1+\alpha)}}\sum_{xy\lambda}\bm r
\end{equation}
\item Convolve $\bm r$ with a smoothing kernel $\bm h$, $\bm r \leftarrow \bm r
\ast \bm h$:
\begin{equation}
    h_{ijk} = \frac{2^{3-|i|-|j|-|k|}}{64} \quad \text{for} \quad i,j,k \in \{-1,0,1\}
\end{equation}
\item Calculate the projections:
\begin{equation}
\tilde {\bm y} = \bm H \bm r
\end{equation}
\item Calculate correction factors for each order $a$, $\bm c^{(a)} \in \mathbb R^{M^2}$:
\begin{equation}
    \bm c^{(a)} = \frac{\bm y^{(a)}}{\tilde {\bm y}^{(a)}}
\end{equation}
\item Back-project each correction factor into the data cube:
\begin{equation}
    \bm C^{(a)} = \bm H^\top_{(a)} \bm c^{(a)} \quad a \in \mathcal S
\end{equation}
where $\bm H_{(a)} \in \mathbb R^{M^2 \times PM^2}$ is the projection operator
for order $a$. For voxels not covered by order $a$'s projection, the correction
is set to the value at the central wavelength slice to prevent boundary
divergence.
\item Update the data cube via the weighted geometric mean of corrections from
all projections that have not yet converged:
\begin{equation}
    \bm r \leftarrow \bm r \prod_{a \in \mathcal S'} \left(\bm C^{(a)}\right)^{w_a/W}
\end{equation}
where $\mathcal S' \subseteq \mathcal S$ is the subset of projections whose
residuals exceed a convergence threshold, $w_a$ is the weight of order $a$, and
$W = \sum_{a \in \mathcal S'} w_a$.
\item Repeat steps 4-7 for $N_\text{inner}$ number of iterations.
\item Repeat steps 2-8 for $N_\text{outer}$ number of iterations.
\end{enumerate}

The algorithm consists of an outer loop between steps 2-8 and an inner loop
between steps 4-7. The outer loop has the additional contrast enhancement
and smoothing steps that can be considered as incorporating the priors for
reconstruction, whereas the inner loop performs corrections to the data cube to
make sure the reconstructed data cube is consistent with the measurements.


The contrast enhancement step serves the purpose of counteracting the unwanted
smearing along the projection direction that the tomographic inversion
algorithms suffer from \citep{kak2001principles}. However, even with the
contrast enhancement, we observe smearing along the spectral dimension of the
data cube, which results in significant overestimation of the line widths. To
overcome that, we impose a spectral prior on the estimated spectra by assuming
that the average of the spectral lines in the data cube is close to the average
two-component spectral template (comprising the primary Fe XII component, the
secondary blend, and the background) derived from the EIS training data (up to
a scaling in intensity). As with the 1D MAP extension, explicitly modeling all
three contributions prevents background emission from being absorbed into the
primary line profile. The spectral prior is implemented as follows. The column-summed spectrum of the
initial cube $\bm r_0$ is computed as $\bm p = \bm S \bm r_0$, where $\bm S$
sums over one spatial dimension, and then scaled column by column to conserve
the total flux of the zeroth-order measurement $\bm y^{(0)}$. This
prior spectrum $\bm p$ is provided to MART as an additional projection with
weight $w_{\rm prior}$, constraining the spectral profiles to a physically
reasonable shape on average. This prior has also been used by
\citet{fox2003data} and \citet{fox2011snapshot} for simulated MOSES inversions,
though they assumed the \textit{actual} $90^\circ$ projection is known, which
is a strong assumption we avoid by using the training data instead.

A limitation of using a fixed average spectral template as the spectral prior is
that it anchors velocity estimates toward the template's mean, suppressing
spatial Doppler structure across the field of view. We address this with a
two-stage procedure. \textit{Stage 1} runs MART with the spectral prior disabled
($w_{\rm prior}=0$) and doubled iteration counts ($2N_\text{outer}$,
$2N_\text{inner}$), using the standard initialization $\bm r_0$. This allows
velocities to evolve freely and produces a spatially-varying Doppler map
$\hat{\bm \mu}^{(1)}$ unconstrained by the average template. However, without
the spectral prior, the recovered line widths incur a large positive bias due
to tomographic smearing. \textit{Stage 2} constructs a velocity-adapted initial cube
\begin{align}
\bm r_0^{(2)} =& \bm \Psi_0(f_1 \bm y^{(0)},\; \hat{\bm \mu}^{(1)},\; \sigma_1^0 \bm 1)
\nonumber\\
              +& \bm \Psi_1(f_2 \bm y^{(0)},\; \hat{\bm \mu}^{(1)} +
              \delta\mu,\; \sigma_2^0 \bm 1) + f_{\rm bg}\, \bm y^{(0)}
\nonumber\\
\end{align}
where $\hat{\bm \mu}^{(1)}$ is the Stage 1 velocity map and
$\delta\mu = (\lambda_2 - \lambda_1)/\Delta\lambda$ is the fixed pixel
separation between the two rest-wavelength centroids. The spectral prior for
Stage 2 is $\bm p^{(2)} = \bm S \bm r_0^{(2)}$, normalized as before. MART is
then re-run from $\bm r_0^{(2)}$ with $w_{\rm prior}=1$, using this
velocity-adapted prior to constrain line widths and intensities while keeping
velocities anchored near the Stage 1 solution.

Once the final data cube is reconstructed, it is fed to the same spectral
fitting routine as we used in the dataset preparation step that solves Eq.
\eqref{eq:spec_fit} to extract the primary line parameters $\bm x$.

\subsection{Comparison Experiments} 
The reconstruction performance of the proposed U-Net is evaluated against the 1D
MAP and MART baselines across varying projection counts $(K)$ and SNR values 
$(\gamma)$. The results reveal clear and consistent performance advantages of
the U-Net over both classical methods for all three reconstructed parameters.
Representative spatial reconstructions are shown in Figure~\ref{fig:comps}.
Figure~\ref{fig:compsweep} provides a visual overview of RMSE and systematic
bias across both experiments, with detailed numerical breakdowns in Tables
\ref{tab:table_k} and \ref{tab:table_gam}. Pixel-level accuracy is examined in
Figure~\ref{fig:scatter}.

\begin{table}[t]
\begin{center}
\setlength{\tabcolsep}{2.8 pt}
\caption{Reconstruction RMSE and Bias for varying projection count $K$ at native
SNR ($\gamma=\infty$). Best performance for each metric is highlighted in bold
(lowest magnitude for Bias). The U-Net achieves the lowest RMSE for all three
parameters at all $K$.}
\vspace{-0.07 in}
\label{tab:table_k}
\begin{tabular}{*8c}
\hline
\textbf{K} & \textbf{Method} & \multicolumn{2}{c}{\textbf{Intensity}} & \multicolumn{2}{c}{\textbf{Velocity}} & \multicolumn{2}{c}{\textbf{Line Width}} \\
\hline
 &  & RMSE & Bias & RMSE & Bias & RMSE & Bias \\
 &  & \multicolumn{2}{c}{\makebox[0pt][c]{(\unit{\erg\per\centi\metre\squared\per\second\per\steradian})}} & \multicolumn{2}{c}{(\unit{\kilo\metre\per\second})} & \multicolumn{2}{c}{(\unit{\kilo\metre\per\second})} \\
\hline\hline
\multirow{3}{*}{2} & U-Net & \textbf{41.9} & -8.5 & \textbf{2.149} & \textbf{-0.250} & \textbf{1.543} & \textbf{-0.048} \\
{} & MART & 56.9 & \textbf{-4.3} & 5.929 & 4.275 & 2.028 & 0.342 \\
{} & \textit{1D MAP} & 58.7 & 18.1 & 5.714 & 1.317 & 4.098 & 0.392 \\
\hline
\multirow{3}{*}{3} & U-Net & \textbf{37.9} & -7.5 & \textbf{1.758} & \textbf{-0.017} & \textbf{1.428} & \textbf{0.072} \\
{} & MART & 54.2 & \textbf{-4.6} & 4.123 & 3.169 & 1.982 & 0.256 \\
{} & \textit{1D MAP} & 58.7 & 18.1 & 5.096 & 1.930 & 3.129 & 0.413 \\
\hline
\multirow{3}{*}{4} & U-Net & \textbf{34.1} & \textbf{-0.8} & \textbf{1.618} & \textbf{-0.042} & \textbf{1.327} & \textbf{0.017} \\
{} & MART & 53.8 & -4.7 & 3.603 & 2.306 & 2.013 & 0.379 \\
{} & \textit{1D MAP} & 58.7 & 18.1 & 4.633 & 1.580 & 3.595 & 1.412 \\
\hline
\multirow{3}{*}{5} & U-Net & \textbf{35.1} & \textbf{-3.7} & \textbf{1.490} & \textbf{0.076} & \textbf{1.233} & \textbf{-0.091} \\
{} & MART & 52.6 & -3.9 & 3.443 & 2.076 & 2.034 & 0.465 \\
{} & \textit{1D MAP} & 58.7 & 18.1 & 4.213 & 1.266 & 3.762 & 1.813 \\
\hline
\end{tabular}
\end{center}
\vspace{-0.20 in}
\end{table}

\begin{table}[t]
\begin{center}
\setlength{\tabcolsep}{2.8 pt}
\caption{Reconstruction RMSE and Bias values of each method evaluated using varying SNR ($\gamma$) levels at $K=3$. The proposed approach maintains the lowest overall errors (RMSE) for all parameters across all noise levels.}
\label{tab:table_gam}
\begin{tabular}{*8c}
\hline
\textbf{$\gamma$} & \textbf{Method} & \multicolumn{2}{c}{\textbf{Intensity}} & \multicolumn{2}{c}{\textbf{Velocity}} & \multicolumn{2}{c}{\textbf{Line Width}} \\
\hline
 &  & RMSE & Bias & RMSE & Bias & RMSE & Bias \\
 &  & \multicolumn{2}{c}{\makebox[0pt][c]{(\unit{\erg\per\centi\metre\squared\per\second\per\steradian})}} & \multicolumn{2}{c}{(\unit{\kilo\metre\per\second})} & \multicolumn{2}{c}{(\unit{\kilo\metre\per\second})} \\
\hline\hline
\multirow{3}{*}{10} & U-Net & \textbf{59.4} & \textbf{-0.9} & \textbf{3.023} & \textbf{0.739} & \textbf{1.821} & \textbf{-0.068} \\
{} & MART & 149.6 & -18.8 & 10.485 & -4.680 & 6.208 & 1.635 \\
{} & \textit{1D MAP} & 139.3 & 18.2 & 9.211 & 2.821 & 2.979 & 2.103 \\
\hline
\multirow{3}{*}{20} & U-Net & \textbf{49.0} & -8.3 & \textbf{2.541} & \textbf{0.221} & \textbf{1.692} & \textbf{-0.140} \\
{} & MART & 87.7 & \textbf{-7.4} & 5.350 & 1.035 & 3.431 & 0.708 \\
{} & \textit{1D MAP} & 87.7 & 18.2 & 8.609 & 2.967 & 2.969 & 1.991 \\
\hline
\multirow{3}{*}{30} & U-Net & \textbf{44.3} & \textbf{-3.6} & \textbf{2.431} & \textbf{0.399} & \textbf{1.619} & \textbf{-0.237} \\
{} & MART & 71.3 & -6.0 & 4.559 & 2.160 & 2.700 & 0.399 \\
{} & \textit{1D MAP} & 73.4 & 18.2 & 7.349 & 2.956 & 3.154 & 1.462 \\
\hline
\multirow{3}{*}{$\infty$} & U-Net & \textbf{37.9} & -7.5 & \textbf{1.758} & \textbf{-0.017} & \textbf{1.428} & \textbf{0.072} \\
{} & MART & 54.2 & \textbf{-4.6} & 4.123 & 3.169 & 1.982 & 0.256 \\
{} & \textit{1D MAP} & 58.7 & 18.1 & 5.096 & 1.930 & 3.129 & 0.413 \\
\hline
\end{tabular}
\end{center}
\end{table}

Across all projection counts, Table \ref{tab:table_k} shows the U-Net achieving
substantially lower RMSE than either classical method in all three parameters.
For velocity at $K=3$, the improvement amounts to roughly a factor of two over
MART and three over 1D MAP, and it holds across the entire range of $K$
evaluated. The line width results require closer inspection. The RMSE values
for the classical methods appear moderate at first glance, but Figure
\ref{fig:scatter} reveals the full picture: MART line width estimates cluster
near the spectral prior value regardless of the true width, and 1D MAP estimates
are broadly scattered. The EIS dataset line width distribution is narrow, with a
standard deviation of only approximately \SI{2}{\kilo\metre\per\second}, so any
method that fails to track spatial variation will achieve modest RMSE simply by
predicting near the mean. The scatter correlation coefficients from Figure
\ref{fig:scatter} confirm this: 0.260 for MART and 0.150 for 1D MAP, compared to
0.724 for the U-Net. Furthermore, the U-Net's velocity and line width accuracy
improve steadily as $K$ increases from 2 to 5, whereas the classical methods
benefit comparatively little from additional projections.

\begin{figure}[h]
\centering
\includegraphics[width=0.5\textwidth]{recons_final.png}
\caption{Ground truth (leftmost column) and reconstructions of each method
for $K=3$ projections at native SNR ($\gamma=\infty$). Rows correspond to
intensity (top), Doppler velocity (middle), and line width (bottom). The U-Net
faithfully recovers the spatial morphology of all three parameters. The 1D MAP
introduces vertical streak artifacts in the velocity and line-width fields due
to its column-wise optimization. MART produces line-width estimates that are
largely insensitive to the true spatial variations.}
\label{fig:comps}
\end{figure}

\begin{figure*}[t]
\centering
\includegraphics[width=1\textwidth]{bar_charts.png}
\caption{Quantitative comparison of reconstruction errors across varying
projection counts ($K$, top row) and signal-to-noise ratios ($\gamma$, bottom
row) for the three physical parameters. To distinguish total error magnitude
from systematic directional error, the charts utilize a dual-encoded visual
design: the wider, semi-transparent bars represent the total root mean square
error (RMSE), while the nested, solid inner bars indicate the mean error (bias).
Across all observational configurations, the proposed U-Net effectively
mitigates the severe systematic overestimations present in the 1D MAP and MART
baselines, consistently yielding the lowest RMSE and near-zero bias.}
\label{fig:compsweep}
\end{figure*}

The U-Net also achieves the lowest intensity RMSE at all $K$ and $\gamma$.
Recovering intensity requires disentangling the dominant spectral line from the
secondary line and coronal background in the zeroth-order projection, combined
with denoising. The classical methods address this by injecting a
multi-component spectral template that encodes the expected contributions of the
primary line, secondary line, and background as prior knowledge, with MART
additionally supplying this template as an explicit spectral projection. Both
also apply spatial smoothness regularization for denoising. The U-Net, by
contrast, learns to extract the primary line parameters directly from the
measurements, without requiring any explicit spectral template. This
data-driven approach also confers a substantial denoising advantage: the
learned representation of coronal spatial structure outperforms simple
smoothness priors at low SNR.

Table \ref{tab:table_gam} quantifies this across the evaluated noise levels.
At the most severe noise level evaluated ($\gamma = 10$), the
U-Net's velocity and line width errors remain below those of either classical
method operating at native noise ($\gamma = \infty$), demonstrating that the
data-driven priors provide a noise robustness that analytical regularizers
cannot match.


\begin{figure*}[t]
\centering
\includegraphics[width=1\textwidth]{joint_scatter.png}
\caption{Hexbin density plots of estimated versus true parameter values across
all 100 test patches at $K=3$ (native SNR), for intensity (top), Doppler
velocity (middle), and line width (bottom). Each hexbin is colored by pixel
count; the Pearson correlation coefficient $r$ is annotated in each panel.
The dashed line indicates the ideal 1:1 correspondence. The U-Net shows the
tightest alignment for all three parameters. MART concentrates line width
estimates near the spectral prior value, reflecting the prior's anchoring
effect on spectral widths. The 1D MAP exhibits the broadest velocity scatter,
consistent with its independent column-wise formulation.}
\label{fig:scatter}
\end{figure*}

Figure \ref{fig:comps} shows representative reconstructions at $K=3$ and
$\gamma=\infty$. While intensity differences between methods are less visually
pronounced, U-Net clearly leads MART and 1D MAP in the morphology of the
velocity and line width reconstructions. The 1D MAP exhibits vertical streak
artifacts in both velocity and line width due to its column-wise optimization.
Neither classical method recovers the line width structure: MART's estimates
are nearly uniform across the field and insensitive to true spatial variations,
while the 1D MAP's reconstruction has high spatial variance that nonetheless
fails to track the true structure. Figure \ref{fig:scatter} confirms these
observations statistically across all 100 test patches.

Once trained, the U-Net reconstructs thousands of patches in seconds, while
the iterative classical methods require hours for the same 100-patch
evaluation set, a practical advantage that complements the accuracy gains
above.



% \subsection{Robustness Analysis}
% \subsubsection{Distribution Shift in Parameter Values}

% \begin{figure}[h]
% \centering
% \includegraphics[width=0.5\textwidth]{paramhists.png}
% % \caption{Histograms of the parameters in the EIS training set}
% \label{fig:paramhist}
% \end{figure}

% \begin{table}[h!]
% \begin{center}
% \caption{Reconstruction SSIM values of each method for the cases of No Shift / Shift}
% \begin{tabular}{|c||c|c|c|}
% \hline
% \textbf{Method} & \textbf{Intensity} & \textbf{Velocity} & \textbf{Line Width} \\
% \hline\hline
% \textit{1D MAP} & 0.920/0.850 & 0.620/0.430 & 0.560/0.410 \\
% \hline
% \textit{SMART} & 0.940/0.910 & 0.570/0.460 & 0.510/0.410 \\
% \hline
% \textit{\textbf{Ours}} & \textbf{0.950/0.930} & \textbf{0.750/0.680} & \textbf{0.710/0.630} \\
% \hline
% \end{tabular}
% \end{center}
% \end{table}

% \subsubsection{Distribution Shift in Spatial Morphologies}

% \subsubsection{Distribution Shift in Signal to Noise}

% \begin{table}[h!]
% \begin{center}
% \caption{Reconstruction SSIM values of each method for the cases of High SNR / Low SNR}
% \begin{tabular}{|c||c|c|c|}
% \hline
% \textbf{Method} & \textbf{Intensity} & \textbf{Velocity} & \textbf{Line Width} \\
% \hline\hline
% \textit{1D MAP} & 0.920/0.850 & 0.620/0.430 & 0.560/0.410 \\
% \hline
% \textit{SMART} & 0.940/0.910 & 0.570/0.460 & 0.510/0.410 \\
% \hline
% \textit{\textbf{Ours}} & \textbf{0.950/0.930} & \textbf{0.750/0.680} & \textbf{0.710/0.630} \\
% \hline
% \end{tabular}
% \end{center}
% \end{table}

\section{Discussion and Conclusions} \label{sec:disc}
We presented a deep learning approach to the spectral inversion problem of
multi-order slitless EUV spectrographs, training a U-Net to map dispersed
observations directly to the intensity, Doppler velocity, and line-width
parameters of the primary coronal emission line. The training data were derived
from a large set of Hinode/EIS rasters, providing a physically realistic
distribution of coronal spectral parameters and spatial morphologies.

Comparison against MART and 1D MAP demonstrated consistent and substantial
improvements across all evaluated configurations. Both classical methods are
fundamentally constrained by the small number of available projections and their
reliance on analytical regularization: MART's smoothness priors cannot resolve
fine spatial structure, while 1D MAP's column-wise optimization introduces
streak artifacts and is sensitive to initialization. By learning spatially
coherent, data-driven priors from the EIS-derived training set, the U-Net
overcomes these constraints. At $K=3$, it reduces velocity RMSE by a factor of
approximately two over MART and three over 1D MAP. The gain in line-width
recovery is more striking: with spatial correlations of 0.260 and 0.150 for
MART and 1D MAP respectively, neither classical method meaningfully tracks true
spatial variation, whereas the U-Net achieves a correlation of 0.724, making
spatially resolved line-width maps a practical output of slitless spectroscopy
for the first time. The learned priors also confer strong noise robustness: at
$\gamma=10$, the U-Net's velocity and line-width errors remain below those of
both classical methods at native SNR. Once trained, inference reduces from hours
to seconds.

These results collectively demonstrate that training-based inversion is a
viable and high-performing approach to spectral recovery from multi-order
slitless EUV observations, and suggest that data-driven methods should be
considered as primary inversion strategies for future solar imaging
spectrographs of this class.

Several limitations of the current work warrant acknowledgment. The methodology
is demonstrated on Hinode/EIS observations of active regions at the
Fe~XII \SI{195.12}{\angstrom} spectral window; applying it to a different
spectral window or plasma regime would require constructing an appropriate
training dataset and retraining the network. Additionally, the current
demonstration uses separately trained networks for each observational
configuration (projection count $K$ and noise level $\gamma$), which limits
flexibility when instrument parameters are not fixed in advance.

Several directions remain open for future work. The generalization properties of
the approach, including the role of training dataset size and sensitivity to
test-time distribution shift, warrant systematic investigation, and improved
network architectures may yield further accuracy gains. Beyond the present
single-line framework, extending the model to accommodate greater spectral
complexity, such as multi-line or non-Gaussian parameterizations, would broaden
its applicability to extreme events and multi-thermal plasmas. A hierarchical
approach, in which a first stage selects the appropriate local spectral model
and a second stage estimates its parameters, provides a natural path toward this
goal. Such extensions, however, introduce a fundamental identifiability
trade-off: because the instrument captures only $K$ projections, there is an
upper bound on the number of free parameters that can be reliably reconstructed.
Incorporating tighter physical constraints through differential emission measure
(DEM) analysis could help overcome this limitation, building upon prior hybrid
tomographic-DEM frameworks \citep{frazin2005combination} and multi-component
decomposition techniques for slitless spectra \citep{cheung2019multi,
winebarger2019unfolding, athiray2025systematic}.

\bibliography{apj25}
\bibliographystyle{aasjournalv7}

%% This command is needed to show the entire author+affiliation list when
%% the collaboration and author truncation commands are used.  It has to
%% go at the end of the manuscript.
%\allauthors

%% Include this line if you are using the \added, \replaced, \deleted
%% commands to see a summary list of all changes at the end of the article.
%\listofchanges

\end{document}

% End of file `sample7.tex'.
